\documentclass{article}
\usepackage{graphicx} % Required for inserting images

\title{STT-2902}
\author{Julien Miron}


\begin{document}

\section*{Équipe 5 – Inférence - Tests d'hypothèses}

Objectif : Expliquer le principe d’inférence statistique à partir d’un ou plusieurs échantillons dans le contexte des tests d'hypothèses.


\vspace{1cm}
Suggestions :
\begin{itemize}
    \item Présentez un scénario où il faut tester une affirmation (ex. : la moyenne d’un groupe est-elle différente d’une valeur cible ?).
    \item Faites réfléchir aux erreurs de type I et II, à la valeur p, et à l’interprétation des résultats du test.
\end{itemize}

\vspace{1cm}
\textbf{Requis} :
\begin{enumerate}
    \item Clairement expliquer la valeur-p.
    \item Illustrer graphiquement les différentes erreurs (type I et II) pour un test donné.
\end{enumerate}

\end{document}
